\documentclass{article}
\usepackage{tcee-repaso}
\tcetema[Política comercial de la UE]{3.B.41}{La política comercial de la Unión Europea.}
\tcefecha{19/11/2026}
\begin{document}
\cabecerarepaso
\begin{columnas}

% ==========================================================================
\bloque[2,5]{Introducción}
\begin{itemize}
\item \textbf{Enganche}: «Donde las dan, las toman».
  \begin{itemize}
  \item 2-4-2025: EE.~UU. anuncia un 20\,\% «recíproco»; la UE aprueba contramedidas por 93.000~M€, pero pacta (27-7-2025) un techo del 15\,\%.
  \item 2026: Mercosur provisional, India y Australia cerradas, México firmado.
  \end{itemize}
\item \textbf{Relevancia}: mayor exportador del mundo; 44~acuerdos con 76~socios = 46,3\,\% de su comercio (2025).
  \begin{itemize}
  \item España: 38,2\,\% de sus exportaciones de bienes, fuera de la UE \vertema{4.A.24}.
  \end{itemize}
\item \textbf{Contextualización}: \concepto{PCC} (arts.~206-207 TFUE): bienes, servicios, PI e IED con terceros; cara exterior de la unión aduanera.
  \begin{itemize}
  \item \emph{Autónoma} (decide sola) y \emph{convencional} (pacta). OMC \vertema{3.B.33}; desarrollo \vertema{3.B.42}.
  \end{itemize}
\item \textbf{Problemática}: ¿por qué común y exclusiva? ¿Cómo se negocia? ¿Qué instrumentos? ¿Cómo con EE.~UU. y China?
\item \textbf{Estructura}: \hyperref[rep:I]{I} bases; \hyperref[rep:II]{II} autónoma; \hyperref[rep:III]{III} convencional y retos.
\end{itemize}
\esquemacante{%
\apartadocante{Introducción}{2,5}
\apartadocante{I. Las bases de la PCC}{7}
\subcante{1. Relevancia}{0,5}
\subcante{2. Teoría económica}{1,5}
\subcante{3. Delimitación de competencias}{1}
\subcante{4. Marco jurídico y decisiones}{2}
\subcante{5. Evolución}{2}
\apartadocante{II. La política comercial autónoma}{9,5}
\subcante{1. La política arancelaria}{2,5}
\subcante{2. Defensa comercial}{3}
\subcante{3. El SPG}{1,5}
\subcante{4. Seguridad económica y sostenibilidad}{2,5}
\apartadocante{III. La política comercial convencional}{10}
\subcante{1. La OMC}{1}
\subcante{2. Acuerdos preferenciales}{5}
\subcante{3. En negociación y pendientes}{1}
\subcante{4. EE.~UU., China y Rusia}{2}
\subcante{5. Retos de futuro}{1}
\apartadocante{Conclusión}{1}}

% ==========================================================================
\bloque[7]{I. Las bases de la política comercial común\label{rep:I}}
\begin{ideaclave}
La PCC nace de la unión aduanera y convierte a 27~economías medianas en una potencia negociadora; su eficacia depende de cómo se repartan el voto y la delegación.
\end{ideaclave}

\sub{I.1. Relevancia\label{rep:I.1}}
\begin{itemize}
\item Extra-UE (2025): bienes 2,65~billones € exportados y 2,54~importados; servicios 1,58~y 1,43.
\item EE.~UU., 1.\textsuperscript{er} socio (21\,\% de las exportaciones; +197.000 M€ en bienes, $-$178.000~en servicios).
\item China, 1.\textsuperscript{er} proveedor ($-$375.000~M€).
\end{itemize}
\graficorepaso[0.95\linewidth]{comercio-socios}{Comercio de bienes con los 16~primeros socios, 2025 (mm€)}

\sub{I.2. Aproximación de teoría económica\label{rep:I.2}}
\subsub{a) ¿Por qué una política comercial?}
\begin{itemize}
\item Librecambio \vertema{3.B.5} \vertema{3.B.6}, salvo arancel óptimo, industria naciente, estratégica \vertema{3.B.7} \vertema{3.B.8}, deslealtad, seguridad.
\item Integración = segundo óptimo (\autor{Viner}, 1950) \vertema{3.B.10}.
\end{itemize}
\subsub{b) ¿Por qué una competencia exclusiva?}
\begin{itemize}
\item Unión aduanera (sin ella, desviación del tráfico); costes de transacción; poder de negociación.
\item Externalidades grandes y preferencias parecidas \vertema{4.B.14}.
\end{itemize}
\subsub{c) La economía política de la «sola voz» (\autor{Meunier}, 2005)}
\begin{itemize}
\item No basta hablar con una voz: decide el \emph{Estado pivote}.
  \begin{itemize}
  \item Con unanimidad, el más reacio: negociador duro en lo defensivo.
  \item Con mayoría cualificada, uno intermedio: acuerdos ofensivos.
  \end{itemize}
\item \autor{Schelling}: manos atadas; \autor{Putnam}: juego a 2~niveles.
\item Explica los «cielos abiertos» y la división del acuerdo con el Mercosur.
\end{itemize}

\sub{I.3. Delimitación de competencias\label{rep:I.3}}
\begin{itemize}
\item \emph{Art.~3.1.e TFUE}: exclusiva, con la unión aduanera (3.1.a) y los acuerdos del art.~3.2.
\item Ámbito (art.~207.1): bienes, servicios, PI comercial e IED.
\item \emph{No es PCC}: promoción (ICEX, CESCE, FIEM) \vertema{4.A.24}; PAC \vertema{3.B.38}; sanciones (PESC, art.~215).
\end{itemize}

\sub{I.4. Marco jurídico y toma de decisiones\label{rep:I.4}}
\subsub{a) El dictamen 2/15 y la nueva arquitectura}
\begin{itemize}
\item \textbf{Dictamen 2/15} (16-5-2017, Singapur): todo exclusivo salvo \emph{inversión de cartera} y \emph{solución inversor-Estado}.
\item \textbf{Arquitectura de 2018}: ALC \emph{solo UE} + acuerdo de inversiones mixto.
  \begin{itemize}
  \item Acuerdos interinos (Chile, México, Mercosur). CETA: 18/27 ratificaciones.
  \end{itemize}
\end{itemize}
\subsub{b) Cómo se negocia y se celebra un acuerdo (arts.~207~y 218)}
\begin{enumerate}
\item La Comisión recomienda; el Consejo autoriza y da directrices.
\item La Comisión negocia con el \concepto{Comité de Política Comercial} e informa al PE.
\item El Consejo firma; el PE aprueba (sin enmiendas); el Consejo celebra (+ 27~si es mixto).
\end{enumerate}
\begin{itemize}
\item Mayoría cualificada; unanimidad en servicios culturales, sociales, educativos y sanitarios (art.~207.4).
\end{itemize}
\graficorepaso[0.95\linewidth]{procedimiento-acuerdos}{Cómo se negocia y se celebra un acuerdo}
\subsub{c) Las reglas de la OMC como marco}
\begin{itemize}
\item Art.~XXIV: lo esencial del comercio, sin subir la protección (Croacia 2013: ajos para China).
\item Cláusula de habilitación: SPG.
\end{itemize}

\sub{I.5. La evolución: de Roma a la seguridad económica\label{rep:I.5}}
\begin{itemize}
\item \emph{Alcance}: Roma (1957); dictamen 1/75 (exclusiva); \textbf{dictamen 1/94} (OMC, mixta); Ámsterdam y Niza (servicios, PI); \textbf{Lisboa} (IED, PLO, aprobación del PE).
\item \emph{Estrategia}:
  \begin{itemize}
  \item multilateralismo $\to$ \emph{Europa Global} (2006: bilaterales);
  \item \emph{Comercio para todos} (2015: valores, transparencia);
  \item revisión de 2021: \concepto{autonomía estratégica abierta} (abierta, sostenible y firme);
  \item \emph{seguridad económica} (2023: promover, proteger, asociarse; doctrina de dic. 2025).
  \end{itemize}
\end{itemize}
\graficorepaso[0.95\linewidth]{cronologia}{La PCC, 1957-2026}

% ==========================================================================
\bloque[9,5]{II. La política comercial autónoma\label{rep:II}}
\begin{ideaclave}
El arancel protege poco salvo en la agricultura; la protección se ha desplazado a la defensa comercial y a instrumentos de coerción, seguridad y sostenibilidad.
\end{ideaclave}

\sub{II.1. La política arancelaria\label{rep:II.1}}
\subsub{a) El arancel aduanero común}
\begin{itemize}
\item \concepto{AAC} (1-7-1968, media de 4~zonas).
\item NMF medio (2024): 5,0\,\% simple (agrícolas 10,5; resto 4,1); 2,8\,\% ponderado.
\item \emph{Simple o ponderada}: la ponderada \emph{infravalora la protección}: lo que tiene arancel alto se importa poco.
  \begin{itemize}
  \item 1988-2022: simple $\approx$6,2 $\to$ 5,2\,\%; ponderada 4,7 $\to$ 2,7\,\%; con preferencias, 1,3\,\%.
  \item Pico de 1995: arancelización de la Ronda Uruguay.
  \end{itemize}
\item Crestas: lácteos 27,1\,\%, bebidas, azúcar, carne; automóviles 10\,\%.
\item Recursos propios (25\,\% para recaudar) \vertema{3.B.37}.
\end{itemize}
\graficorepaso[0.75\linewidth]{aranceles-ue}{Arancel NMF medio de la UE por productos, 2024 (\%)}
\graficorepaso[0.9\linewidth]{aranceles-evolucion}{Arancel medio de la UE, 1988-2022: simple y ponderado (\%)}
\subsub{b) La nomenclatura: del SA al TARIC}
\begin{itemize}
\item 2~dígitos, capítulo; 4, partida; 6, subpartida SA (mundial); 8, NC (derecho de terceros); 10, TARIC (medidas).
\item Bicicleta 8712~00~30: 14\,\%.
  \begin{itemize}
  \item China: +48,5\,\% antidumping; Indonesia: SPG 10,5\,\% pero antidumping extendido; Vietnam: 0\,\% (ALC).
  \end{itemize}
\end{itemize}
\subsub{c) Modulación y regímenes aduaneros}
\begin{itemize}
\item Contingentes, suspensiones, franquicias; depósito, zona franca, perfeccionamiento.
\item Reforma 2026: Código 2026/2108, Autoridad Aduanera (Lille), fin de la franquicia de 150 € (3 € por artículo).
\end{itemize}

\sub{II.2. Las medidas de defensa comercial\label{rep:II.2}}
\subsub{a) Antidumping (Reg. 2016/1036)}
\begin{itemize}
\item Margen $\geq$2\,\%, perjuicio, causalidad e \emph{interés de la Unión}.
\item Regla del derecho inferior (no si hay distorsiones en insumos, 2018); \emph{distorsiones significativas} (2017/2321).
\item Plazos: denuncia 25\,\%; provisionales 7-8 meses; definitivas 14; 5~años; compromisos de precios.
\item Adopta la Comisión (actos de ejecución).
\end{itemize}
\subsub{b) Antisubvenciones y salvaguardias}
\begin{itemize}
\item \emph{Antisubvenciones} (2016/1037): subvención específica.
  \begin{itemize}
  \item VE chinos (2024): 7,8-35,3\,\% + 10\,\%; China a la OMC (DS630) y represalias (brandy, porcino, lácteos).
  \end{itemize}
\item \emph{Salvaguardias}: acero 2018-30-6-2026 (8~años máx.).
  \begin{itemize}
  \item Desde 1-7-2026, Reg. 2026/1384 (art.~XXVIII): cupos 18,3~Mt ($-$47\,\%), 50\,\% fuera, \emph{melt and pour}.
  \end{itemize}
\item Obstáculos al comercio (2015/1843); \emph{Enforcement} (654/2014, reformado en 2021).
\end{itemize}
\subsub{c) La defensa comercial en cifras}
\begin{itemize}
\item 232~medidas en vigor (fin de 2025; 199~un año antes); 32~investigaciones en 2025 (media $\approx$12).
\end{itemize}
\graficorepaso[0.92\linewidth]{defensa-ue}{Investigaciones iniciadas y medidas impuestas por la UE, 1995-2025}
\subsub{d) ¿A quién protege el antidumping? (\autor{Prusa}, 1992)}
\begin{itemize}
\item \emph{Colusión inducida}: el compromiso de precios coordina; la empresa europea sube su precio.
  \begin{itemize}
  \item Ganan las 2, a costa de consumidores e industrias usuarias.
  \end{itemize}
\item \emph{Evidencia}: importaciones investigadas $-$30/50\,\%, con desviación (\autor{Prusa}, 2001); menos en la UE (\autor{Konings} y otros, 2001).
\item Por eso, el \emph{interés de la Unión}.
\end{itemize}

\sub{II.3. El Sistema de Preferencias Generalizadas\label{rep:II.3}}
\begin{itemize}
\item UNCTAD (1968); CEE (1971); cláusula de habilitación (1979).
\item 3~regímenes:
  \begin{itemize}
  \item \emph{general}: renta baja y media-baja, graduación;
  \item \emph{SPG+}: 8~países que aplican convenios (Bolivia, Filipinas, Pakistán, Sri Lanka…);
  \item \emph{TMA} (2001): PMA.
  \end{itemize}
\item Retiradas: Myanmar, Bielorrusia, Camboya (2020).
\item Reg. 2026/1395 (2027-2036): condicionalidad de \emph{readmisión}; sale quien firma un ALC (Indonesia, Kenia).
\end{itemize}

\sub{II.4. Seguridad económica y sostenibilidad\label{rep:II.4}}
\begin{itemize}
\item \emph{Seguridad}:
  \begin{itemize}
  \item \textbf{Anticoerción} (2023/2675): Consejo por mayoría cualificada; nunca activado: disuasión;
  \item \textbf{IPI} (2022/1031): productos sanitarios chinos $>$5~M€ (2025);
  \item subvenciones extranjeras (2022/2560); control de inversiones (2026/1386) \vertema{3.B.39}; doble uso (2021/821).
  \end{itemize}
\item \emph{Sostenibilidad}:
  \begin{itemize}
  \item \textbf{MAFC} (2023/956): definitivo desde 1-1-2026; certificados desde 2-2027 \vertema{3.B.31};
  \item \textbf{deforestación}: 30-12-2026; \textbf{diligencia debida}: Ómnibus I (2026/470), 2029;
  \item \textbf{trabajo forzoso}: 12-2027.
  \end{itemize}
\end{itemize}

% ==========================================================================
\bloque[10]{III. La política comercial convencional\label{rep:III}}
\begin{ideaclave}
Ante el bloqueo de la OMC, la UE ha tejido la mayor red de acuerdos (53,3\,\% del comercio con los ya cerrados); sus 2~primeros socios quedan fuera y con ellos manda el poder de negociación.
\end{ideaclave}

\sub{III.1. La OMC\label{rep:III.1}}
\begin{itemize}
\item Examen cada 3~años; ACP (22~partes), ATI, aeronaves; MPIA (2020, 62~miembros) \vertema{3.B.33} \vertema{3.B.34}.
\end{itemize}

\sub{III.2. Los acuerdos comerciales preferenciales\label{rep:III.2}}
\subsub{a) Tipos y generaciones}
\begin{enumerate}
\item Vecindad y descolonización: Turquía, Yaundé-Lomé.
\item Mercancías: EEE (1994), Sudáfrica, México (2000), Chile (2003).
\item Nueva generación: Corea (2011), Colombia-Perú y Centroamérica (2013), CETA (2017), Japón y Singapur (2019), Vietnam (2020).
\item Autonomía estratégica: N. Zelanda (2024), Kenia (2024), Chile interino (2025).
\end{enumerate}
\begin{itemize}
\item Capítulos de desarrollo sostenible con sanciones (revisión de 2022).
\end{itemize}
\graficorepaso[0.85\linewidth]{escalera-preferencias}{La escalera de las preferencias de la UE}
\subsub{b) La red de acuerdos en vigor}
\begin{itemize}
\item \emph{Europa}: EEE (1994); Suiza (ALC 1972, 120~bilaterales, paquete de 2026); uniones aduaneras (Turquía, Andorra, San Marino); Balcanes (AEA); DCFTA (Ucrania, Moldavia, Georgia).
\item \emph{Reino Unido}: ACC (2021, solo UE): sin aranceles ni cuotas con reglas de origen; igualdad de condiciones; sin pasaporte.
\item \emph{América}:
  \begin{itemize}
  \item CETA (2017, provisional; 99\,\% de líneas; 143~IG);
  \item Mercosur (interino, 1-5-2026); México (2000; modernizado, firmado 2026);
  \item Chile (interino 2025: litio, cobre); Colombia-Perú-Ecuador y Centroamérica (2013).
  \end{itemize}
\item \emph{Asia-Oceanía}: Corea (2011); Japón (2019; 200~IG); Singapur (2019); Vietnam (2020); N. Zelanda (2024).
\item \emph{ACP}: AAE asimétricos; euromediterráneos.
\end{itemize}
\graficorepaso[0.95\linewidth]{mapa-acuerdos}{La red de acuerdos de la UE, noviembre de 2026}
\subsub{c) ¿Qué efecto tienen los acuerdos?}
\begin{itemize}
\item 46,3\,\% del comercio (2025); con socios preferenciales, +3,1\,\% frente a +1,4\,\%.
\item CETA: +71\,\% (2016-2023); un ALC duplica el comercio en 10~años (\autor{Baier}-\autor{Bergstrand}, 2007).
\end{itemize}

\sub{III.3. Acuerdos en negociación y pendientes\label{rep:III.3}}
\begin{itemize}
\item \textbf{Mercosur}: cerrado 12-2024; firmado 17-1-2026; interino provisional 1-5-2026 \vertema{3.B.35}.
  \begin{itemize}
  \item Asociación, pendiente del dictamen 1/26; liberaliza 91/92\,\%; vacuno 99.000~t al 7,5\,\%; 344~IG.
  \end{itemize}
\item México firmado (22-5-2026); India cerrado (27-1-2026); Indonesia (9-2025); Australia (24-3-2026).
\item En negociación: Tailandia, Filipinas, Malasia, EAU; China (CAI), congelado (2021).
\end{itemize}

\sub{III.4. Estados Unidos, China y Rusia\label{rep:III.4}}
\subsub{a) Estados Unidos}
\begin{itemize}
\item TTIP parado (2019); acero 2018; Airbus-Boeing (tregua 2021).
\item 2025: acero 50\,\%, autos 25\,\%; Turnberry (27-7) y declaración (21-8): techo del 15\,\% todo incluido.
  \begin{itemize}
  \item UE a cero en industriales; 750.000~M\$ en energía, 600.000~M\$ de inversión.
  \end{itemize}
\item El Supremo anula los IEEPA (20-2-2026) $\to$ s.~122~y s.~301~con techo; contramedidas suspendidas.
\end{itemize}
\subsubsub{Negociación bajo amenaza (\autor{Nash}, 1950)}
\begin{itemize}
\item Sin reglas, el reparto lo decide el \emph{punto de desacuerdo}: gana quien menos pierde con la ruptura (\autor{Grossman}-\autor{Helpman}, 1995).
\item Explica el 15\,\% frente a 0~y el valor disuasorio de las contramedidas y del Anticoerción (RU 10\,\%; UE y Japón 15\,\%).
\end{itemize}
\subsub{b) China, Rusia y ACP}
\begin{itemize}
\item \textbf{China}: socio, competidor y rival sistémico (2019); \emph{de-risking}; tierras raras (2025).
\item \textbf{Rusia}: 21~paquetes de sanciones (7-2026); gas fuera en 2027; del 5.\textsuperscript{o} socio (2021) al 19.\textsuperscript{o} (2025).
\item ACP: Acuerdo de Samoa \vertema{3.B.42}.
\end{itemize}

\sub{III.5. Retos de futuro\label{rep:III.5}}
\begin{itemize}
\item Multilateralismo sin la OMC de antes (MPIA, plurilaterales).
\item Estabilidad con EE.~UU. (declaración sin solución de diferencias).
\item Exceso de capacidad chino sin caer en el proteccionismo.
\item Ratificar a tiempo (CETA, Mercosur, India, Indonesia, Australia).
\item Sostenibilidad frente a competitividad (MAFC, deforestación, Ómnibus).
\end{itemize}

% ==========================================================================
\bloque[1]{Conclusión}
\begin{itemize}
\item \textbf{Recapitulación}: unión aduanera $\to$ poder negociador; autónoma (arancel, defensa, SPG, seguridad) y convencional (la mayor red de acuerdos).
\item \textbf{Relevancia}: España: acceso a mercados (aceite, vino, automóvil) y defensa \vertema{4.A.24}.
\item \textbf{Extensiones}: \vertema{3.B.33} \vertema{3.B.10} \vertema{3.B.35} \vertema{3.B.39} \vertema{3.B.42}.
\item \textbf{Opinión}: la seguridad económica no debe acabar protegiendo a quien más presiona.
  \begin{itemize}
  \item La mejor defensa, la red de acuerdos, por mayoría cualificada.
  \end{itemize}
\item \textbf{Idea final}: la amenaza solo pesa si es creíble y si hay alternativas; una potencia comercial no se defiende cerrándose.
\end{itemize}

% ==========================================================================
\anexorepaso
\bloque{Datos clave}
\begin{itemize}
\item Extra-UE: \dato{bienes 2.648/2.542 mm€}{2025, exp./imp.}.
\item Arancel NMF: \dato{5,0\,\% simple; 2,8\,\% ponderado}{OMC, 2024}.
\item Red: \dato{44 acuerdos, 46,3\,\% del comercio}{2025}; 53,3\,\% con los cerrados.
\item Defensa: \dato{232 medidas}{fin de 2025}.
\item España: \dato{38,2\,\% de las exportaciones de bienes, fuera de la UE}{2025}.
\end{itemize}
\bloque{Autores y fechas}
\hito{1950}{\autor{Viner}; \autor{Nash}: negociación}
\hito{1968}{Arancel aduanero común}
\hito{1992}{\autor{Prusa}: colusión inducida}
\hito{1994}{Dictamen 1/94}
\hito{2005}{\autor{Meunier}: la sola voz}
\hito{2009}{Lisboa: IED en la PCC}
\hito{2017}{Dictamen 2/15}
\hito{2025}{Turnberry: techo del 15\,\%}
\bloque{Preguntas del tribunal}
\begin{enumerate}
\item \textbf{¿Por qué es exclusiva la PCC?} Unión aduanera, costes de transacción, poder de negociación.
\item \textbf{¿Qué dijo el dictamen 2/15?} Exclusiva salvo inversión de cartera y solución inversor-Estado.
\item \textbf{¿Qué es el interés de la Unión?} Prueba propia de la UE: usuarios y consumidores.
\item \textbf{¿Por qué aceptó la UE el 15\,\%?} Peor punto de desacuerdo (superávit en bienes, seguridad).
\item \textbf{¿Por qué la ponderada es menor que la simple?} Lo que tiene arancel alto se importa poco.
\item \textbf{¿Se ha usado el Anticoerción?} No; su fin es disuadir.
\end{enumerate}
\begin{trampa}
\begin{itemize}
\item La IED entra en la PCC con \emph{Lisboa} (2009).
\item El antidumping \emph{no} lo aprueba el Consejo desde 2014.
\item Turquía: unión aduanera solo \emph{industrial}.
\item Suiza \emph{no} está en el EEE; el ACC con el RU es \emph{solo UE}.
\item El techo del 15\,\% \emph{incluye} el NMF; el acero sigue al 50\,\%.
\end{itemize}
\end{trampa}

\end{columnas}
\end{document}
